\documentclass[11pt]{article}
\usepackage{explorer}
\newcommand{\ExplorerLanguage}{FR}
\newcommand{\ExplorerNumber}{12}
\newcommand{\ExplorerKind}{Solutions détaillées}
\begin{document}
\ExplorerTitle{Terminale · analyse et algorithmique}{Encadrer une racine par dichotomie}
\ExplorerSection{Données des exercices}
Pour l’exercice 1 : on cherche un zéro de \(f(x)=x^{2}-2\) sur [1 ; 2]. Après deux partages, l’intervalle conservé est [1,25 ; 1,5].\par
\ExplorerSection{Exercice 1}
Poursuivre l’exemple une troisième fois.\par
Le milieu vaut 1,375 ; son carré vaut \(1,890625<2\). Le nouvel intervalle est [1,375 ; 1,5].\par
\ExplorerSection{Exercice 2}
Partant d’un intervalle de largeur 1, combien de partages garantissent une largeur au plus \(\frac{1}{1000}\) ?\par
\(2^{9}=512<1000\leq 1024=2^{10}\), donc 10 partages suffisent et 9 ne suffisent pas.\par
\ExplorerSection{Exercice 3}
Pourquoi le changement de signe de \(f(x)=\frac{1}{x}\) entre −1 et 1 ne prouve-t-il pas l’existence d’un zéro ?\par
La fonction n’est pas définie en 0 et n’est pas continue sur [−1,1]. L’hypothèse du théorème n’est pas satisfaite ; \(\frac{1}{x}\) n’est jamais nul.\par
\ExplorerSection{Exercice 4}
Une racine est dans [1,4 ; 1,5]. Quelle erreur maximale garantit le milieu 1,45 ?\par
La demi-largeur vaut \(\frac{1,5-1,4}{2}=0,05\).\par
\ExplorerSection{Problème : raisonnement complet}
Pour \(f(x)=x^{3}-x-1\) sur [1,2], prouver l’existence et l’unicité d’une racine \(\alpha \). Effectuer trois partages par dichotomie. Combien de partages garantissent que le milieu final approche \(\alpha \) à 0,001 près, avec une erreur au plus 0,001 ?\par
\(f(1)=-1\) et \(f(2)=5\) ; f est continue. De plus \(f'(x)=3x^{2}-1\geq 2\) sur [1,2], donc f est strictement croissante. Les milieux 1,5, 1,25, 1,375 donnent successivement [1 ; 1,5], [1,25 ; 1,5], [1,25 ; 1,375] car les valeurs de f sont \(0,875, -0,296875\) et 0,224609375. Après n partages, l’erreur du milieu est au plus \(\frac{1}{2^{n+1}}\). Il faut \(2^{n+1}\geq 1000\), donc \(n\geq 9\).\par
\end{document}
